\documentclass[UTF8,fontset=none,10pt,a4paper]{ctexart}

\usepackage[a4paper,margin=1.2cm]{geometry}
\usepackage{fontspec}
\usepackage{xeCJK}
\usepackage{booktabs}
\usepackage{array}
\usepackage{enumitem}
\usepackage{listings}
\usepackage{xcolor}
\usepackage{titlesec}
\usepackage{hyperref}

\hypersetup{
  colorlinks=true,
  linkcolor=black,
  urlcolor=blue
}

\setmainfont{TeX Gyre Termes}
\setsansfont{TeX Gyre Heros}
\setmonofont{Latin Modern Mono}[Scale=0.90]
\setCJKmainfont{Noto Serif CJK TC}
\setCJKsansfont{Noto Sans CJK TC}
\setCJKmonofont{Noto Sans CJK TC}[Scale=0.90]
\xeCJKsetup{
  CJKecglue=\hspace{0.12em},
  CheckSingle=true
}

\setlength{\parindent}{1.8em}
\setlength{\parskip}{0.18em}
\linespread{1.0}
\setlist[itemize]{leftmargin=1.4em,itemsep=0.05em,topsep=0.1em}
\setlist[enumerate]{leftmargin=1.4em,itemsep=0.05em,topsep=0.1em}
\titleformat{\section}{\large\bfseries\sffamily}{\thesection}{0.55em}{}
\titlespacing*{\section}{0pt}{0.55em}{0.2em}
\titleformat{\subsection}{\normalsize\bfseries}{\thesubsection}{0.5em}{}
\titlespacing*{\subsection}{0pt}{0.35em}{0.15em}

\lstset{
  basicstyle=\ttfamily\scriptsize,
  columns=fullflexible,
  breaklines=true,
  frame=single,
  framerule=0.2pt,
  xleftmargin=0.5em,
  xrightmargin=0.5em,
  aboveskip=0.3em,
  belowskip=0.3em
}

\newcommand{\code}[1]{\texttt{\detokenize{#1}}}

\title{\vspace{-1.0em}\sffamily\bfseries AI-Based AIG Optimization Final Project Report}
\author{傅豐樺 \quad 學號：114501702}
\date{\vspace{-0.6em}}

\begin{document}
\maketitle
\vspace{-1.2em}
\fontsize{9.2pt}{10.6pt}\selectfont

\section{專案目標與重現方式}

本專案針對 \code{benchmarks/ex200.truth} 到 \code{ex299.truth} 共 100 個 Boolean function，產生功能等價且 area-delay product (ADP) 較低的 AIG。ADP 定義為 $\mathrm{area}\times\mathrm{delay}$，其中 area 是 AIG 中 two-input AND node 數量，delay 是 AIG level。所有輸出都必須通過等價檢查，否則不列入結果。

主要程式與工具放在 \code{student/}。其中 \code{frontends/} 負責 truth-table 解析、Verilog simulation 與 Yosys synthesis，\code{generators/} 產生不同 benchmark family 的 Verilog/RTL，\code{backends/} 執行 AIG 最佳化，\code{scripts/} 負責整理 best result 與更新 \code{output/}。重現目前提交用的 \code{output/exNNN.aig} 可在專案根目錄執行：

\begin{lstlisting}
cd ALS_Final_Project
export PATH=$PWD/student/tools/conda-env/bin:$PATH
python3 student/scripts/update_current_best_by_case.py
python3 student/scripts/update_output_best.py
python3 evaluate.py
\end{lstlisting}

\code{update_current_best_by_case.py} 會掃描已驗證結果並選出每個 case 最低 ADP 的 AIG。\code{update_output_best.py} 會複製成 \code{output/ex200.aig} 到 \code{output/ex299.aig}，並在 \code{student/results/output_current_best_index.csv} 記錄來源、area、delay、ADP 與 hash。

這份 index 也是本報告採用的實驗依據。它不是額外人工整理的表，而是從目前 \code{output/} 實際使用的 AIG 回推來源。每一列都包含輸出檔案、原始 run、來源 CSV、AIG hash 與複製時間，因此可以檢查某個 case 的結果到底來自 frontend seed、後端 GA 還是後續 rescue run。這樣做的好處是報告中的總 ADP 與 case-level 分析都對應到最後提交的 \code{output/}，而不是只對應某次中途實驗。

前端解的來源是 AI-based AIG generation。流程是先由 AI 輔助拆解 truth table 背後的語意，得到 high-level description，再用 generator 或 AI-assisted Verilog rewriting 產生多種 Verilog 寫法，最後交給 Yosys 合成 AIG，並用 \code{evaluate.py} 量測 area、delay、ADP 與等價性：

\begin{center}
\small
\code{truth table -> AI reverse engineering -> high-level description -> Verilog generators -> Yosys -> AIG -> evaluate.py}
\end{center}

若要回到更前面的 reverse engineering 階段，則從對應 case 的 \code{source_run}、\code{source_csv} 或 \code{student/frontend_campaigns/} index 找到當時使用的 generator 與 hypothesis，再重跑該 generator 產生 Verilog。若要重新跑某個 domain 的後端 GA，可使用：

\begin{lstlisting}
python3 student/scripts/run_backend_portfolio.py \
  --portfolio-version v3-ga --domain integer --cases ex255-ex279 \
  --run-id newseeds_integer_v3a_20260614_1952 \
  --current-bundle student/new_seeds/integer/ex255_ex279_current \
  --backend-bundle-glob 'student/new_seeds/integer/ex255_ex279_backend_candidates_*' \
  --jobs 40 --case-jobs 10 --task-jobs 4 \
  --skip-culs --disable-esyn \
  --max-generations 8 --population-size 32 \
  --offspring-per-generation 96 --max-sequence-len 4 \
  --per-case-runtime-soft-limit 7200 --abc-timeout 90 --resume
\end{lstlisting}

其他 benchmark family 使用同一套指令結構，只需替換 \code{--domain}、\code{--cases}、\code{--run-id} 與候選 AIG 路徑。最後一律以 \code{python3 evaluate.py} 作為正式等價性與 ADP 檢查。

重現時我沒有要求所有中間搜尋都重新跑完，因為部分 GA 或 deep synthesis run 需要數小時以上。報告中的「可重現」主要分成兩層：第一層是從已驗證結果重建 \code{output/} 並跑 \code{evaluate.py}，這是最終提交必須通過的檢查。第二層才是回到 frontend campaign 或 backend run 的 source path，重跑當時的 generator 或 command sequence。前者確保最後答案可驗證，後者則保留方法來源與除錯線索。

在實際整理結果時，我也盡量讓每個輸出都能追溯到明確來源。除了 \code{output_current_best_index.csv} 之外，相關 run 的 \code{best.csv}、\code{candidates.csv} 與工作目錄會保留 candidate id、來源 AIG、command sequence、area/delay 與 hash。這些欄位讓我可以分辨某個改善是來自新的 frontend Verilog、Yosys synthesis flow，還是後端 GA 的某段指令序列。對報告來說，這比只列最後分數更重要，因為它能說明方法是否真的可重做。

\section{方法架構}

整體 flow 分成前端候選產生、候選整理、後端 GA、結果彙整四層：

\begin{center}
\small
\begin{tabular}{@{}p{0.22\linewidth}p{0.72\linewidth}@{}}
\toprule
階段 & 內容 \\
\midrule
前端候選產生 & 從 truth table 反推語意或結構化 Verilog，再經 Yosys/ABC 合成為 AIG。\\
候選整理 & 將已驗證且表現較好的 AIG 按 case 整理，作為後端搜尋起點。\\
後端 GA & 用 genetic algorithm 在多種 synthesis command sequence 中搜尋更低 ADP。\\
結果彙整 & 收斂所有通過驗證的結果，更新 \code{current_best_by_case.csv} 與 \code{output/}。\\
\bottomrule
\end{tabular}
\end{center}

前端候選產生的核心是避免直接展開 truth table。直接展開通常會得到很大的 AIG，因此本專案先用 AI 做 reverse engineering，把原本看起來像隨機 bit pattern 的 truth table 轉成 high-level description，再寫成結構化 Verilog。BF16/FP16 主要拆 sign/exponent/mantissa 與 special cases，FP8/float 類 cases 使用 field-class decision graph、alignment 與 residual overlay，integer cases 使用 multiplier/divider/square/isqrt shell，unknown cases 則嘗試 BDD、ANF/Davio、cofactor、prefix/key transform 或 image/route descriptor。每個 Verilog candidate 都經 Yosys/ABC synthesis 成 AIG，再用 \code{evaluate.py} 檢查等價性與 ADP。

前端搜尋時，我通常不只相信單一 semantic guess。即使兩段 Verilog 在功能上等價，Yosys/ABC 合成後的 AIG 結構可能差很多，所以同一個 high-level description 會被改寫成多種 RTL：例如使用 case split、lookup table、decision tree、arithmetic expression 或 shared predicate。這些寫法會造成不同的 area/delay tradeoff，也會影響後端 GA 是否容易再改善。因此前端階段留下多個 verified candidates，比只保留當下最低 ADP 的 candidate 更有價值。

候選整理時，我會把每個 case 分成「目前最好」與「可作為後端起點」兩類。前者負責提供目前最低 ADP 的基準，後者則保留結構多樣性。舉例來說，有些 AIG 的 ADP 略差，但 delay 較低或邏輯結構較平衡，後端 GA 反而可能從這類起點找到更好的結果。因此 \code{student/new_seeds/} 的功能不是代表前端方法本身，而是把已驗證的前端成果整理成後端搜尋可使用的 seed set。

後端 GA 不再修改 truth table 或 Verilog，而是在已等價的 AIG 上搜尋 optimization command sequence。GA 的 individual 可以視為「起始 AIG 加上一串 synthesis 指令」。每個 case 會建立自己的 population，初始化時放入目前最佳 AIG、其他等價候選，以及過去接近最佳的結果。每一代從較好的 result 中選 parent，產生 offspring。產生方式包括改動 command sequence、重組兩個 parent 的指令片段，或從另一個起始 AIG 重新搜尋。保留結果時不只看最低 ADP，也保留不同 area/delay tradeoff，避免搜尋太早集中在單一路線。

後端可選指令主要來自 ABC rewrite/refactor/resub/balance/dc2，以及 ABC9、MockTurtle 的等價重寫與平衡。每個 offspring 執行後都重新量測，只有通過 \code{evaluate.py} 且 ADP 有效的 AIG 會被保留。不等價、timeout 或明顯退化的結果會被丟棄。整體來看，後端 GA 把大量可能的 synthesis script 當成搜尋空間，用實際驗證後的 ADP 作為回饋訊號，逐步改善每個 case。

後端 GA 的另一個重點是保留 diversity。若只保留 ADP 最低的一個 AIG，後續搜尋常會很快收斂到同一種結構，改善幅度有限。因此我在整理候選時會保留不同 root family、不同 delay、不同 area 的結果，讓 population 可以從多個方向出發。對接近 reference 的 case，搜尋會偏向短 sequence 和低風險 rewrite。對落後較多的 case，才會嘗試更長 sequence、較重的 synthesis flow 或從另一個 frontend seed 重新開始。

GA 的 fitness 以 ADP 為主，但我沒有完全忽略 area 與 delay 的個別變化。原因是 ADP 相同時，低 delay 的 AIG 和低 area 的 AIG 在後續 rewrite 中可能有不同潛力。實作上會記錄每個 candidate 的 area、delay、ADP、hash 與來源指令，並用 timeout 控制單一任務成本。若中途停止，後續可以用 resume 繼續接著跑，避免已花時間產生的候選被浪費。

\section{AI/LLM 在流程中的使用}

AI/LLM 在本專案中主要用在實驗設計、假設生成、程式實作與結果分析，並不直接決定最終 AIG 是否正確。正確性仍由 ABC 等價檢查與 \code{evaluate.py} 決定。

第一，LLM 協助做 benchmark family classification 與 high-level function hypothesis。它根據 truth-table pattern、case notes 與既有 run summary，把 cases 分成 BF16、FP16、FP8/float、integer、unknown 等類別，再提出可能的 decomposition。第二，LLM 用來快速建立與修改工具，例如 truth-table analyzer、Verilog generator、後端 GA wrapper、結果彙整與 output update scripts。這些工具讓實驗結果可以固定輸出成 CSV 與 manifest，後續才有辦法重現和比較。

第三，LLM 也被用來平行化前端探索。對於語意不明的 cases，我把不同 case 或不同 hypothesis 分給多個 sub-agent 同時處理。有些 sub-agent 專看 BF16/FP16 的 exponent/mantissa pattern，有些專攻 integer multiplier/divider/isqrt，有些則針對 unknown family 嘗試 BDD、cofactor、ANF 或 image/route descriptor。每個 sub-agent 都要產生可驗證的 Verilog、Yosys synthesis 結果與 evaluation CSV，而不是只回報文字猜測。最後再收斂各 agent shard 的成果與失敗假設，留下真正等價且 ADP 有競爭力的候選。

我對 sub-agent 的要求是產物必須能被主流程驗證。也就是說，文字推論本身不算結果，至少要能留下 generator、Verilog、synthesis log 或 evaluation row。若 hypothesis 失敗，也會記錄失敗原因，例如不等價、ADP 過高、只改善單一 output、或合成後結構爆炸。這些失敗紀錄能避免下一輪重複嘗試同一條路，也讓後續 prompt 可以更精準地限制搜尋方向。

另外，我也用 LLM 把外部方法快速轉成可實驗版本。E-Syn 的核心概念是用 e-graph 保存多種等價 Boolean expression，再抽出 area 或 delay 較好的表示。我根據其 paper 與 GitHub 實作建立了 \code{student/backends/esyn_flow.py} 與 \code{student/backends/esyn_eqn.py}，將 AIG 轉成 EQN/S-expression 後做 per-output rewrite，再接回 AIG 並用 \code{evaluate.py} 檢查。這個版本最後可以作為 optional seed generator，但實測改善不明顯，因此正式結果沒有依賴 E-Syn。

整體來說，LLM 在這個專案裡比較像加速實驗迭代的工具，而不是 correctness oracle。它可以幫忙整理 hypotheses、產生腳本、根據失敗結果提出下一輪方向，也可以讓多個 sub-agent 同時探索不同 family。但所有 AI 產生的 Verilog、AIG 或 command sequence 最後都要通過同一套等價檢查與 ADP 量測。這個界線很重要，因為 truth table benchmark 很容易出現看起來合理但其實不等價的猜測。

\section{實驗結果與觀察}

本節結果以 \code{student/results/output_current_best_index.csv} 為準。這份檔案對應目前實際放在 \code{output/} 的 100 個 AIG，欄位包含每個 case 的 area、delay、ADP、reference ADP、來源 run、來源 CSV、AIG hash 與複製時間。也就是說，下面的統計不是中途某次 run 的摘要，而是目前提交輸出的結果。此版本於 2026-06-15 12:11:46 +0800 更新到 \code{output/}。

目前 100 個 case 都有輸出 AIG，且皆可由 \code{evaluate.py} 做正式檢查。總 ADP 為 7,462,959，reference 總 ADP 為 6,696,028，總和 ratio 約 1.115。若以逐 case 的 \code{ADP/reference ADP} 來看，45 個 case 已優於 reference，68 個 case 低於 1.1x reference，96 個 case 低於 1.5x reference。平均 ratio 約 1.083，median ratio 約 1.012，表示大部分 case 已接近 reference，總分主要被少數落後較多的 case 拉高。

表格中的 \emph{Sum ratio} 是該 domain 的 $\sum ADP / \sum reference\ ADP$，因此會反映大型 case 對總分的影響。\emph{Avg ratio} 則是逐 case ratio 的平均，比較能看出該 domain 中一般 case 的表現。這兩個數字有時會不同，例如 BF16 的 Sum ratio 低於 1，但 Avg ratio 略高於 1，代表少數 ADP 較大的 case 改善很多，而其他 case 大多只是接近 reference。

\begin{center}
\small
\begin{tabular}{lrrrrr}
\toprule
Domain & Cases & Sum ratio & Avg ratio & Beat ref & $>$1.5x \\
\midrule
BF16 & 20 & 0.888 & 1.034 & 10 & 0 \\
FP16 & 20 & 1.058 & 1.057 & 10 & 1 \\
FP8/float & 15 & 0.569 & 0.840 & 10 & 0 \\
Integer & 25 & 1.073 & 1.075 & 9 & 0 \\
Unknown & 20 & 1.369 & 1.351 & 6 & 3 \\
\bottomrule
\end{tabular}
\end{center}

從 domain summary 來看，FP8/float 是目前效果最明顯的一組，總和 ratio 只有 0.569，15 個 case 中有 10 個 beat reference。這組的改善主要來自前端語意重建，因為這些 function 若能被解讀成 field-level operation 或較小的 decision graph，Yosys/ABC 合成後會比直接展開 truth table 精簡許多。BF16 的總和 ratio 也低於 1，但平均 ratio 略高於 1，代表有幾個大 case 改善明顯，整體分布則仍接近 reference。Integer family 的結果較平均，沒有特別嚴重的 outlier，但也不像 FP8/float 那樣有大幅勝出的 case。

相對地，Unknown family 是目前主要瓶頸。雖然 \code{ex287} ratio 只有 0.488，顯示部分 unknown case 可以透過 decomposition 找到更好的結構，但 \code{ex286}、\code{ex297}、\code{ex299} 仍明顯落後 reference。這幾個 case 的共同問題是前端結構還不夠明確，後端 GA 即使能做局部 rewrite，也很難把一個不夠好的初始結構完全修回來。因此後續若要繼續改善總分，重點應該放在這些 high-ratio case 的 frontend reconstruction，而不是只增加一般後端搜尋時間。

\begin{center}
\small
\begin{tabular}{lrrp{0.53\linewidth}}
\toprule
Case & Ratio & ADP & 觀察 \\
\midrule
\code{ex251} & 0.187 & 7,344 & frontend semantic seed 明顯優於 reference，是目前改善最大的 case。\\
\code{ex254} & 0.414 & 15,840 & FP8/float 類型受益於較好的 high-level structure。\\
\code{ex287} & 0.488 & 2,820 & unknown 中少數被 decomposition 打開的 case，顯示此類方法仍有潛力。\\
\code{ex286} & 5.947 & 14,131 & 結構仍不明確，單靠後端 command sequence 改善有限。\\
\code{ex297} & 1.764 & 398,400 & ADP 本身很大，仍需要更好的 frontend guess 或 shared predicate。\\
\bottomrule
\end{tabular}
\end{center}

整體而言，本次結果的主要收穫是：前端找對語意時，後端 GA 可以進一步把 area/delay 壓低，甚至大幅 beat reference；但若前端只得到偏展開式或不夠貼近語意的 AIG，後端搜尋很容易進入局部最佳。這也解釋了為什麼結果中同時存在 ratio 低於 0.2 的成功案例，以及少數 ratio 高於 1.5 的瓶頸案例。

從提交角度看，目前結果已經涵蓋全部 100 個 benchmark，且所有輸出都來自 \code{source_kind=run} 的 verified rows。這代表最後的 \code{output/} 不是手動拼接的黑箱，而是由已記錄來源的 run 結果整理而來。後續若要更新成更好的版本，只要新的 run 產生更低 ADP 且通過等價檢查，就可以重新執行 update scripts 產生新的 index 與 output。

\section{困難、限制與改進方向}

第一個困難是搜尋空間很大。即使只考慮 ABC/ABC9 的 command sequence，不同起始 AIG、指令順序、timeout 與 area/delay tradeoff 都會造成結果差異。本專案透過保留多種 tradeoff、記錄指令過去成效、支援 resume，以及設定每個 case 的時間上限來控制成本，但仍不能保證找到 global optimum。

第二個困難是前端初始 AIG 的品質常比後端最佳化更重要。有些 case 透過後端 GA 可以把 ratio 從高於 1.5 拉回 1.1 附近，但也有 case 反覆優化仍停在同一區域，代表 AIG 結構本身不夠好。這類 case 需要更好的 semantic reconstruction，例如 unknown family 的 image/route descriptor、cofactor key、class one-hot mux，或 FP16 selected-bit/exception shell。

第三個限制是部分工具受環境影響。需要 GPU 的最佳化 flow 在本地環境中不一定穩定。E-Syn 目前只保守用在單一輸出 expression rewrite，避免一次替換所有輸出造成轉換成本過高。耗時的深度 synthesis flow 也只適合針對少數落後較多的 cases 做最後補強，而不適合放入每個 case 的一般搜尋。

未來可以改進的方向包括：建立學習式 command selection，讓後端 GA 利用 case feature 預測哪些操作可能有效。對 unknown functions 建立更系統化的 truth-table feature extractor。最後，針對落後較多的 cases 做更深入的 frontend decomposition，而不是只增加後端時間。

另外，報告中的結果也顯示「前端語意」與「後端最佳化」其實不是兩個獨立階段。好的 frontend seed 會決定後端能不能找到更好的局部區域，後端搜尋結果也會反過來提醒哪些 semantic guess 可能不對。未來若有更多時間，我會把這兩者串成更自動化的 loop：當後端長時間無法改善時，自動回到 frontend 產生不同 decomposition，而不是只延長同一條 backend search。


\end{document}
